\chapter{La digestión y los nutrientes}\label{ch:g7:digestion-and-nutrients}

El \cref{ch:g4:journey-of-food} siguió el mordisco de pan tubo abajo y
prometió que ``los \omterm{def:g4:journey-of-food:digestion}{jugos digestivos} lo desmontan''. Este curso abrimos
las puertas del taller: ¿qué hacen exactamente los jugos, cuáles son
los \omterm{def:g7:digestion-and-nutrients:families}{nutrientes} que producen y cómo se las apaña la pared del \omterm{def:g4:journey-of-food:tract}{intestino
delgado} con el mayor paso fronterizo del cuerpo? El mordisco de pan
está de vuelta --- con mejores instrumentos.

\section{Los jugos en faena}

\begin{proposition}[La digestión es un desmontaje químico]\label{prop:g7:digestion-and-nutrients:chemical}
Masticar y revolver solo rompen el alimento en trozos \emph{más
pequeños}; el desmontaje de verdad es químico. Los \emph{jugos
digestivos}\index{jugo digestivo} --- la \omterm{def:g4:journey-of-food:tract}{saliva}, el jugo del \omterm{def:g4:journey-of-food:tract}{estómago},
los jugos que se vierten al \omterm{def:g4:journey-of-food:tract}{intestino delgado} --- contienen sustancias
de trabajo que cortan los componentes grandes del alimento en otros
pequeños: el almidón en azúcares, la materia de construcción de la
carne y de los \omterm{def:g2:animals-grow-up:born}{huevos} en sus unidades pequeñas, las grasas en las
suyas. Cada jugo tiene sus especialidades y cada uno funciona mejor en
las \omterm{def:g6:exploring-our-environment:conditions}{condiciones} de su propia estación.
\end{proposition}
\admitted

\begin{example}[La versión de laboratorio]\label{ex:g7:digestion-and-nutrients:bench}
El experimento clásico hace la \omterm{def:g4:journey-of-food:digestion}{digestión} dentro del vidrio. Dos tubos
con papilla de almidón cocido a \omterm{def:g6:exploring-our-environment:conditions}{temperatura} del cuerpo: a uno se le
echa \omterm{def:g4:journey-of-food:tract}{saliva} fresca y al otro \omterm{def:g4:journey-of-food:tract}{saliva} hervida (estropeada). Una hora
después, la prueba del almidón encuentra el tubo 2 sin cambios ---
sigue siendo almidón --- mientras que la papilla del tubo 1 se ha
convertido en azúcares, exactamente el dulzor del
\cref{ex:g4:journey-of-food:sweet}, producido sin ninguna boca. El jugo
solo hizo el desmontaje --- y el testigo hervido enseña que su
sustancia de trabajo la destruye el calor, como le pasaría a una
maquinaria delicada.
\end{example}

\begin{remark}[¿Y para qué desmontar?]\label{rem:g7:digestion-and-nutrients:why}
Los componentes grandes del alimento pueden atravesar la pared del
intestino tan poco como un armario puede pasar por una ranura de
buzón. Y el cuerpo no los quiere enteros: quiere \emph{piezas} ---
unidades pequeñas que pueda reconstruir como sustancia propia
(\cref{prop:g6:matter-from-food:transform}). La \omterm{def:g4:journey-of-food:digestion}{digestión} es un
desmontaje para el transporte \emph{y} para el montaje: muebles
convertidos en kit desmontable, en cada comida.
\end{remark}

\section{Los nutrientes}

\begin{definition}[Las familias de nutrientes]\label{def:g7:digestion-and-nutrients:families}
El desmontaje da los \emph{nutrientes}\index{nutriente} --- los
componentes pequeños y listos para la \omterm{prop:g4:journey-of-food:absorb}{sangre}:
\begin{itemize}
  \item los \emph{azúcares} (sobre todo la \omterm{prop:g7:organs-at-work:bill}{glucosa}, el combustible de
        la \cref{prop:g7:organs-at-work:bill}), de los alimentos
        feculentos y dulces;
  \item las unidades pequeñas de la \emph{materia de construcción}, de
        la carne, el pescado, los \omterm{def:g2:animals-grow-up:born}{huevos}, los lácteos y las legumbres
        --- los ladrillos de la construcción de la
        \cref{prop:g3:balanced-diet:jobs};
  \item las unidades de las \emph{grasas} --- a la vez combustible
        denso y material de construcción;
  \item más unos pasajeros que no necesitan desmontaje: las
        \emph{vitaminas}, las \emph{\omterm{prop:g4:what-plants-need:needs}{sales minerales}} y la propia agua.
\end{itemize}
Las \omterm{def:g3:balanced-diet:families}{familias de alimentos} de la \cref{def:g3:balanced-diet:families}
son, en el fondo, proveedoras de estas familias de nutrientes --- la
lista de la compra que hay detrás de la lista de la compra.
\end{definition}

\begin{example}[Una comida, desglosada]\label{ex:g7:digestion-and-nutrients:lunch}
La comida equilibrada del \cref{ex:g3:balanced-diet:lunch}, releída
como \omterm{def:g7:digestion-and-nutrients:families}{nutrientes}: la pasta --- almidón, desmontado en \omterm{prop:g7:organs-at-work:bill}{glucosa} a litros;
el pescado --- materia de construcción, cortada en sus unidades; el
aceite de las judías --- unidades de grasa; las judías y la manzana ---
vitaminas, minerales, agua y además fibra. Fíjate en que la fibra es la
excepción que confirma el sistema: ninguno de nuestros jugos la puede
cortar, así que no atraviesa nada y sigue su camino --- el volumen de
las sobras (\cref{prop:g4:journey-of-food:absorb}), haciendo su propio
barrido útil.
\end{example}

\section{El gran paso fronterizo}

\begin{proposition}[La absorción en las vellosidades]\label{prop:g7:digestion-and-nutrients:absorption}
El paso que prometía la \cref{prop:g4:journey-of-food:absorb} ocurre a
través de la pared interior del \omterm{def:g4:journey-of-food:tract}{intestino delgado}, y la pared está
hecha para eso: plegada y alfombrada con millones de
\emph{vellosidades}\index{vellosidad} --- salientes en forma de dedo de
un milímetro de alto, cada uno envolviendo sus propios \omterm{def:g4:heart-and-blood:vessels}{vasos
sanguíneos} finos. Las \omterm{prop:g7:digestion-and-nutrients:absorption}{vellosidades} son la razón de que la superficie de
paso llegue a su tamaño de alfombra de una habitación
(\cref{rem:g4:journey-of-food:surface}): pliegues sobre pliegues sobre
dedos. Los \omterm{def:g7:digestion-and-nutrients:families}{nutrientes} atraviesan sus paredes finas hasta la \omterm{prop:g4:journey-of-food:absorb}{sangre}, y
la \omterm{prop:g4:journey-of-food:absorb}{sangre} los lleva primero a un gran órgano que clasifica y almacena
--- el \emph{hígado}\index{hígado} --- antes del reparto general.
\end{proposition}
\admitted

\begin{omfigure}
\begin{tikzpicture}[scale=0.9]
  % intestine wall with villi
  \draw[thick] (0,0) -- (0,3.4);
  \draw[thick] (0,0) -- (8.2,0);
  \node[font=\small, rotate=90] at (-0.4,1.7) {pared del intestino};
  % villi
  \foreach \x in {0.9,2.4,3.9,5.4,6.9} {
    \draw[thick, fill=omDef!8] (\x,0)
      .. controls ({\x-0.42},1.4) and ({\x-0.42},2.4) .. (\x,2.6)
      .. controls ({\x+0.42},2.4) and ({\x+0.42},1.4) .. (\x,0);
    % vessels inside
    \draw[omThm, thick] ({\x-0.14},0.15)
      .. controls ({\x-0.18},1.3) .. (\x,2.25)
      .. controls ({\x+0.18},1.3) .. ({\x+0.14},0.15);
  }
  \node[font=\footnotesize, align=center] at (4.1,-0.6)
    {cada vellosidad envuelve sus propios vasos sanguíneos finos};
  % nutrients crossing
  \foreach \p in {(1.7,2.9),(3.2,3.0),(4.7,2.85),(6.2,2.95)} {
    \fill[omMeth] \p circle (0.07);
  }
  \node[font=\footnotesize, omMeth] at (6.9,3.25)
    {nutrientes en el tubo};
  \draw[->, omMeth, thick] (3.2,2.9) -- (3.65,2.2);
  \node[font=\footnotesize, align=left] at (10.0,1.7)
    {el paso:\\ por la pared fina\\ de la vellosidad,\\ hasta la sangre};
\end{tikzpicture}

{\small La alfombra interior del intestino: millones de \omterm{prop:g7:digestion-and-nutrients:absorption}{vellosidades},
cada una un dedo de pared fina lleno de \omterm{def:g4:heart-and-blood:vessels}{vasos sanguíneos}. Superficie
plegada sobre superficie --- el gran paso fronterizo del cuerpo.}
\end{omfigure}

\begin{example}[Otra vez el pliego de condiciones]\label{ex:g7:digestion-and-nutrients:spec}
Grandes, finas, húmedas y forradas de \omterm{prop:g4:journey-of-food:absorb}{sangre}: las \omterm{prop:g7:digestion-and-nutrients:absorption}{vellosidades} cumplen
exactamente el pliego de \omterm{def:g6:exploring-our-environment:conditions}{condiciones} de la
\cref{prop:g7:how-animals-breathe:spec} --- para \omterm{def:g7:digestion-and-nutrients:families}{nutrientes} en vez de
para \omterm{def:g7:what-is-respiration:respiration}{oxígeno}. Bolsas de los \omterm{def:g7:how-animals-breathe:four}{pulmones}, \omterm{def:g1:animals-around-us:covering}{plumas} de las \omterm{def:g7:how-animals-breathe:four}{branquias},
\omterm{prop:g7:digestion-and-nutrients:absorption}{vellosidades} del intestino: allí donde el cuerpo tiene que mover mucho
a través de una frontera, construye la misma máquina. Reconoce el
pliego una vez y lo reconocerás en todas partes.
\end{example}

\begin{method}[Rastrear un nutriente de punta a punta]\label{met:g7:digestion-and-nutrients:trace}
Para cualquier bocado, rastrea un \omterm{def:g7:digestion-and-nutrients:families}{nutriente}:
\begin{enumerate}
  \item nombra la familia del alimento y sus componentes principales;
  \item estación por estación, anota dónde empieza y dónde acaba su
        desmontaje y qué jugos han trabajado en él;
  \item atraviesa con él las \omterm{prop:g7:digestion-and-nutrients:absorption}{vellosidades}, hasta la \omterm{prop:g4:journey-of-food:absorb}{sangre}, pasando por
        el \omterm{prop:g7:digestion-and-nutrients:absorption}{hígado};
  \item entrégalo: a una \omterm{def:g5:first-look-at-cells:cell}{célula} que trabaja, como combustible
        (\cref{prop:g7:organs-at-work:bill}), o a una obra, como
        material (\cref{def:g6:matter-from-food:production}).
\end{enumerate}
\end{method}

\begin{example}[El rastreo, hecho con pan]\label{ex:g7:digestion-and-nutrients:breadtrace}
Pan: familia feculenta, sobre todo almidón. El desmontaje empieza en la
boca (la \omterm{def:g4:journey-of-food:tract}{saliva} --- el dulzor de masticar mucho rato), hace una pausa
en el ácido del \omterm{def:g4:journey-of-food:tract}{estómago} y termina en los jugos del \omterm{def:g4:journey-of-food:tract}{intestino delgado}:
\omterm{prop:g7:organs-at-work:bill}{glucosa}. Paso: \omterm{prop:g7:digestion-and-nutrients:absorption}{vellosidades}, \omterm{prop:g4:journey-of-food:absorb}{sangre}, \omterm{prop:g7:digestion-and-nutrients:absorption}{hígado} --- que ingresa parte de la
\omterm{prop:g7:organs-at-work:bill}{glucosa} para más tarde, uno de sus muchos servicios. Entrega: un
\omterm{def:g3:bones-and-muscles:muscle}{músculo} del muslo la quema a media hora de esprint. El mordisco del
cuarto curso, con las cuentas completas.
\end{example}

\section{Ejercicios}

\begin{exercise}[$\star$]\label{exo:g7:digestion-and-nutrients:1}
¿Qué hacen los \omterm{def:g4:journey-of-food:digestion}{jugos digestivos} que no pueden hacer los dientes ni el
revuelto?
\end{exercise}

\begin{exercise}[$\star$]\label{exo:g7:digestion-and-nutrients:2}
Describe el experimento de los dos tubos con \omterm{def:g4:journey-of-food:tract}{saliva}: montaje, resultado
y las dos conclusiones.
\end{exercise}

\begin{exercise}[$\star$]\label{exo:g7:digestion-and-nutrients:3}
Enumera las familias de \omterm{def:g7:digestion-and-nutrients:families}{nutrientes} y una fuente alimentaria de cada una.
\end{exercise}

\begin{exercise}[$\star$]\label{exo:g7:digestion-and-nutrients:4}
¿Qué componentes de los alimentos no necesitan desmontaje? ¿Y cuál
derrota a nuestros jugos por completo --- y para qué fin útil?
\end{exercise}

\begin{exercise}[$\star$]\label{exo:g7:digestion-and-nutrients:5}
¿Qué es una vellosidad? ¿Qué envuelve cada una?
\end{exercise}

\begin{exercise}[$\star$]\label{exo:g7:digestion-and-nutrients:6}
¿Adónde va la \omterm{prop:g4:journey-of-food:absorb}{sangre} que viene del intestino antes del reparto general
y qué servicio presta ese órgano? Nombra uno.
\end{exercise}

\begin{exercise}[$\star\star$]\label{exo:g7:digestion-and-nutrients:7}
Explica la imagen del kit desmontable de la
\cref{rem:g7:digestion-and-nutrients:why}: las dos razones por las que
hay que desmontar el alimento.
\end{exercise}

\begin{exercise}[$\star\star$]\label{exo:g7:digestion-and-nutrients:8}
¿Por qué importa el tubo de la \omterm{def:g4:journey-of-food:tract}{saliva} hervida? Nombra el método y qué
enseña el hervido sobre la sustancia de trabajo del jugo.
\end{exercise}

\begin{exercise}[$\star\star$]\label{exo:g7:digestion-and-nutrients:9}
Pon las \omterm{prop:g7:digestion-and-nutrients:absorption}{vellosidades} al lado de las bolsas de los \omterm{def:g7:how-animals-breathe:four}{pulmones} y de las
\omterm{def:g1:animals-around-us:covering}{plumas} de las \omterm{def:g7:how-animals-breathe:four}{branquias}: enuncia el pliego de \omterm{def:g6:exploring-our-environment:conditions}{condiciones} compartido y
la carga de cada máquina.
\end{exercise}

\begin{exercise}[$\star\star$]\label{exo:g7:digestion-and-nutrients:10}
Pásale el \cref{met:g7:digestion-and-nutrients:trace} a un \omterm{def:g2:animals-grow-up:born}{huevo} cocido
(materia de construcción), hasta una obra de reparación en un rasguño
que se cura (\cref{exo:g6:cell-unit-of-life:10}).
\end{exercise}

\begin{exercise}[$\star\star$]\label{exo:g7:digestion-and-nutrients:11}
Una dieta radical se salta todos los alimentos de materia de
construcción durante semanas. Con la
\cref{def:g7:digestion-and-nutrients:families} y la
\cref{def:g6:matter-from-food:production}, predice qué se queda corto y
qué sufre primero en un cuerpo que todavía crece.
\end{exercise}

\begin{exercise}[$\star\star\star$]\label{exo:g7:digestion-and-nutrients:12}
``La \omterm{def:g4:journey-of-food:digestion}{digestión} convierte una comida en una lista de reparto''. Defiende
la frase de punta a punta --- jugos, \omterm{def:g7:digestion-and-nutrients:families}{nutrientes}, \omterm{prop:g7:digestion-and-nutrients:absorption}{vellosidades}, \omterm{prop:g7:digestion-and-nutrients:absorption}{hígado},
entregas --- en un párrafo.
\end{exercise}

\section{Problema: El tubo digestivo de vidrio}

\begin{problem}[{Problema de fin de semana --- construir la digestión
en el laboratorio, tubo a tubo}]\label{pb:g7:digestion-and-nutrients:1}
La clase construye un ``\omterm{def:g4:journey-of-food:tract}{tubo digestivo} de vidrio'': una gradilla de
tubos a \omterm{def:g6:exploring-our-environment:conditions}{temperatura} del cuerpo, cada uno simulando una estación del
\omterm{def:g4:journey-of-food:tract}{tubo digestivo}. Tubo M (boca): papilla de almidón más \omterm{def:g4:journey-of-food:tract}{saliva} fresca.
Tubo M$'$: papilla de almidón más \omterm{def:g4:journey-of-food:tract}{saliva} hervida. Tubo S (\omterm{def:g4:journey-of-food:tract}{estómago}):
cubitos de clara de \omterm{def:g2:animals-grow-up:born}{huevo} más jugo de \omterm{def:g4:journey-of-food:tract}{estómago} de la farmacia. Tubo
S$'$: esos mismos cubitos en agua sola. Todos funcionan dos horas; y
después, las pruebas --- del almidón, de los azúcares --- y un vistazo
a los cubitos.

\medskip\noindent\textbf{Parte I --- Predicciones y resultados.}
\begin{enumerate}
  \item Predice los resultados a las dos horas de los cuatro tubos.
  \item M se vuelve azucarado y M$'$ no. Enuncia la conclusión y el
        papel exacto de M$'$.
  \item Los cubitos de S encogen, se ablandan y se aclaran; los de
        S$'$ se quedan como estaban. ¿Qué ha hecho el jugo de \omterm{def:g4:journey-of-food:tract}{estómago}
        --- y qué clase de familia de \omterm{def:g7:digestion-and-nutrients:families}{nutrientes} se está produciendo?
  \item Ningún tubo se masticó ni se revolvió. ¿Qué demuestra la
        gradilla sobre dónde está el trabajo de verdad de la \omterm{def:g4:journey-of-food:digestion}{digestión}?
\end{enumerate}

\medskip\noindent\textbf{Parte II --- Mejorar el modelo.}
\begin{enumerate}[resume]
  \item Un alumno propone hacer funcionar el tubo M a \omterm{def:g6:exploring-our-environment:conditions}{temperatura} de
        nevera ``por orden''. Objeta, con el patrón de la
        \cref{prop:g6:decomposers-and-soil:speed} aplicado a los jugos.
  \item Otro propone un solo tubo grande con los dos jugos mezclados.
        ¿Qué hecho de estación por estación de la
        \cref{prop:g7:digestion-and-nutrients:chemical} emborronaría la
        mezcla?
  \item El \omterm{def:g4:journey-of-food:tract}{tubo digestivo} de vidrio no tiene pared: no se absorbe nada.
        ¿Qué estructura necesitaría a continuación un modelo completo y
        qué dos propiedades tiene que tener su material?
  \item Añade fibra --- copos de salvado --- a un quinto tubo con todos
        los jugos disponibles. Predice su estado a las dos horas y qué
        significa para el modelo.
\end{enumerate}

\medskip\noindent\textbf{Parte III --- Del vidrio al intestino.}
\begin{enumerate}[resume]
  \item En el cuerpo, los azúcares de M se encontrarían a continuación
        con las \omterm{prop:g7:digestion-and-nutrients:absorption}{vellosidades}. Continúa su viaje en cuatro paradas.
  \item El jugo del \omterm{def:g4:journey-of-food:tract}{estómago} trabaja en un ácido fuerte; los jugos del
        intestino necesitan que ese ácido esté neutralizado. ¿Qué dice
        esto sobre lo de ``cada jugo en las \omterm{def:g6:exploring-our-environment:conditions}{condiciones} de su propia
        estación'' --- y sobre tragarse las estaciones en orden?
  \item Una empresa farmacéutica vende cápsulas de jugos para personas
        cuyas propias estaciones rinden poco. Explica, con la gradilla,
        por qué una terapia así es siquiera posible.
  \item Resume la práctica en una frase que contenga \emph{jugos},
        \emph{\omterm{def:g7:digestion-and-nutrients:families}{nutrientes}} y \emph{tubos testigo}.
\end{enumerate}
\end{problem}
